% Owner: Member A (Kashaf Ali); mapping assumed from proposal order.
\chapter{Introduction}
This chapter is mandatory. The introduction should contain a brief overview of the problem being addressed and the background information needed for the reader to understand the work being done and the reasoning behind it.  
The last paragraph of the introduction chapter should contain an outline of the entire report.  Summarize each chapter in one line to make the last paragraph.
\section{Purpose of this Document}
\textbf{Specify the purpose of this Project.}
The purpose of the document in an FYP report is to clearly and concisely state the main objective and goals of the project, as well as the research question(s) or hypothesis(es) being investigated. It should provide a brief overview of the background and context of the project, and explain why the research is important and relevant. The purpose of the document should be written in a way that is easy for the reader to understand, and should be included at the beginning of the report.\\ In addition to providing an overview of the project, the purpose of the document should also outline the specific objectives and goals of the research, such as what the project aims to achieve and what the research questions or hypotheses are. This will help the reader understand the scope of the project and what to expect from the report. It is also important to include a brief description of the methods and techniques used in the research, as well as any limitations or assumptions of the project. This will provide the reader with a clear understanding of how the research was conducted and what the results and conclusions are based on. Finally, it should be clear what the report is going to cover, what kind of information will be presented, and in what order.\\
For Example, "The purpose of this document is to present the design, implementation, and evaluation of our Final Year Project (FYP) which aims to develop a web-based application for small businesses to manage their inventory and customer orders. The project will focus on creating a user-friendly interface and providing a seamless user experience. The goal of this research is to provide a Minimum Viable Product (MVP) that can be used by small businesses to improve their inventory management and customer order processing. The research question we aim to answer is: Can we develop a web-based application that can improve inventory management and customer order processing for small businesses? This report will present the methodology, design, implementation, testing, and evaluation of our project, as well as the limitations and future work."
\section{Intended Audience}
Identifying and writing the intended audience for a Final Year Project (FYP) report is an important step in the writing process. The intended audience is the group of people for whom the report is written and who will primarily read and use the report. It is crucial to identify the intended audience early in the writing process as this will guide the report's style, tone, and level of detail.\\
To identify the intended audience, consider the report's purpose, who will be using the report and for what purpose. For example, the intended audience for a FYP report in Computer Science may be a panel of professors who will be evaluating the project. In contrast, the intended audience for a FYP report in Business may be potential investors or clients.
\section{Definitions, Acronyms, and Abbreviations}
List all important definitions, acronyms, and abbreviations used in this document. For example: 
\textbf{SDG}: Sustainable Development Goal\\
\textbf{FYP:} Final Year Project\\
\textbf{MVP}: Minimum Viable Product\\
\textbf{UI}: User Interface\\
\textbf{UX}: User Experience\\
\textbf{Agile}: Agile Development\\
\textbf{SCRUM}: Scrum Development\\
\textbf{REST}: Representational State Transfer\\
\textbf{ORM}: Object-Relational Mapping\\
\textbf{CRUD}: Create, Read, Update, Delete\\
\textbf{API}: Application Programming Interface
\section{Conclusion}
The last paragraph of the introduction chapter should contain an outline of the entire report. Summarize each chapter in one line to make the last paragraph.
